%! TeX program = lualatex
% An example talk using tamu-talk. Copy it as a starting point for your own.
% The Texas A&M logo is not included: see the README or `texdoc tamu-talk`.
\DocumentMetadata{
  lang=en,
  tagging=on,
  pdfstandard=UA-2,
  pdfstandard=A-4f,
}

\documentclass{ltx-talk}

\usepackage{tamu-talk}

\SetKeys[tamu-talk]{
  page-numbers = total,
  title-bold   = false,
  title-box    = true,
  header-style = banner,
  title-logo   = false, % set to true once the logo is installed
}

\usepackage{fontspec}
\setmainfont{Ysabeau}

\newcommand{\lining}[1]{{\addfontfeature{Numbers=Lining}#1}}

\setmonofont{Fira Mono}[Scale=MatchLowercase]

\usepackage{lua-unicode-math}
\usepackage{fancyvrb}

\title{Title of Your Talk}
\subtitle{An Optional Subtitle}
\author{John Doe}
\date{\today}

\begin{document}

\begin{frame}
  \maketitle[
    framestyle    = tamu,
    title-size    = 18pt,
    subtitle-size = 11pt,
    logo-width    = .25\textwidth,
    logo-height   = 28mm,
  ]
\end{frame}

\section{Equation}

\begin{frame}
  \frametitle{An Equation}
  \[
    \int_0^\infty e^{-x^2} \, dx = \frac{\sqrt{\pi}}{2}
  \]
\end{frame}

\section{Figure}

\begin{frame*}
  \frametitle{A Figure with Alt Text}
  Alt text is supplied via the \texttt{alt=} key on \texttt{\textbackslash includegraphics}.

  \bigskip
  \begin{columns}
    \begin{column}{0.55\textwidth}
      \begin{sidebar}
        \begin{Verbatim}[gobble=8]
        \begin{figure}
          \centering
          \includegraphics[
            width=.5\linewidth,
            alt={A gray placeholder image
              with a large letter A.}
          ]{example-image-a}
        \end{figure}
        \end{Verbatim}
      \end{sidebar}
    \end{column}
    \begin{column}{0.4\textwidth}
      \begin{figure}
        \centering
        \includegraphics[
          width=.55\linewidth,
          alt={A gray placeholder image with a large letter A.}
        ]{example-image-a}
      \end{figure}
    \end{column}
  \end{columns}
\end{frame*}

\section{Boxes}

\begin{frame*}
  \frametitle{headerbox}
  \begin{headerbox}{Title}
    Solid maroon header bar, tan-tinted body, one fixed style.
  \end{headerbox}

  \bigskip
  \begin{sidebar}
    \begin{Verbatim}[gobble=6]
      \begin{headerbox}{Title}
        Body text.
      \end{headerbox}
    \end{Verbatim}
  \end{sidebar}
\end{frame*}

\begin{frame*}
  \frametitle{lasso}
  \begin{lasso}{Title}
    Solid border, left-offset floating title tab.
  \end{lasso}

  \bigskip
  \begin{lasso}[black]{Title}
    The only other color option.
  \end{lasso}

  \bigskip
  \begin{sidebar}
    \begin{Verbatim}[gobble=6]
      \begin{lasso}[maroon|black]{Title}
        Body text.
      \end{lasso}
    \end{Verbatim}
  \end{sidebar}
\end{frame*}

\begin{frame*}
  \frametitle{sidebar}
  \begin{sidebar}
    Thick maroon left bar, tan-tinted body, no title.
  \end{sidebar}

  \bigskip
  \begin{sidebar}
    \begin{Verbatim}[gobble=6]
      \begin{sidebar}[variant] % (optional variants: forest | slate)
        Body text.
      \end{sidebar}

      % ...or any tcolorbox keys, e.g. to recolor it by hand
      \begin{sidebar}[colback=TAMUGray!20, borderline west={3pt}{0pt}{TAMUGray}]
        Body text.
      \end{sidebar}
    \end{Verbatim}
  \end{sidebar}
\end{frame*}

\begin{frame*}
  \frametitle{elevate}
  \begin{elevate}[Title]
    Rounded corners, maroon header bar, drop shadow.
  \end{elevate}

  \bigskip
  \begin{elevate}
    No title given -- no header bar at all.
  \end{elevate}

  \bigskip
  \begin{sidebar}
    \begin{Verbatim}[gobble=6]
      \begin{elevate}[Title]  % [Title] is optional; omit it for no header
        Body text.
      \end{elevate}
    \end{Verbatim}
  \end{sidebar}
\end{frame*}

\begin{frame*}
  \frametitle{outline}
  \begin{outline}
    No fill at all -- just a thin maroon border.
  \end{outline}

  \bigskip
  \begin{outline}[2pt]
    A custom line width.
  \end{outline}

  \bigskip
  \begin{sidebar}
    \begin{Verbatim}[gobble=6]
      \begin{outline}[.75pt]  % line width, optional, default shown
        Body text.
      \end{outline}
    \end{Verbatim}
  \end{sidebar}
\end{frame*}

\section{Colors}

\newcommand{\swatch}[2]{%
  \begin{minipage}[t]{3.2cm}
    {\fboxsep=0pt \colorbox{#1}{\rule{0pt}{1.2cm}\rule{3.2cm}{0pt}}}
    \\[2pt]
    \small #1\par
    \lining{#2}
  \end{minipage}%
}

\begin{frame}
  \frametitle{Brand Colors}

  \begin{center}
    \swatch{TAMUMaroon}{\#500000}
    \swatch{TAMUTan}{\#D6D3C4}
    \swatch{TAMUGray}{\#707070}
  \end{center}

  \bigskip
  \begin{center}
    \swatch{TAMUDarkGray}{\#3E3E3E}
    \swatch{TAMULightGray}{\#D1D1D1}
    \swatch{TAMUNearBlack}{\#202020}
  \end{center}
\end{frame}

\begin{frame*}
  \frametitle{Non-Brand Colors}
  The package \verb|ltx-talk| itself provides an \verb|\alert{}| command which uses \alert{this red} color by default.
  \vspace{1em}

  It also provides a named color called \verb|example| which is \textcolor{example}{green}.
  \vspace{1em}

  Additionally, \verb|tamu-talk| provides two off-brand colors that pair well with maroon:

  \bigskip
  \begin{sidebar}[forest]
    \textbf{forest} (off-brand).
  \end{sidebar}

  \bigskip
  \begin{sidebar}[slate]
    \textbf{slate} (off-brand).
  \end{sidebar}
\end{frame*}

\appendix

\section{Appendix}

\subsection{Backup}

\begin{frame*}
  \frametitle{Supplementary Material}
  Use \verb|\appendix| to section off backup slides (to avoid running up the page counter).

  \bigskip
  % gobble=N strips the first N characters of leading whitespace from
  % every line of the Verbatim block. It lets this block be indented in
  % the .tex source (to read cleanly alongside the surrounding code)
  % without that indentation showing up in the typeset output.
  \begin{sidebar}
    \begin{Verbatim}[gobble=6]
      \begin{document}

      % ... main talk content ...

      \appendix

      \begin{frame}
        \frametitle{Supplementary Material}
        Body text.
      \end{frame}

      \end{document}
    \end{Verbatim}
  \end{sidebar}
\end{frame*}

\subsection{Overlays}

\begin{frame}
  \frametitle{Incremental Overlays}
  \begin{itemize}[<+->]
    \item First point, shown on slide 1 onward.
    \item Second point, shown on slide 2 onward.
    \item Third point, shown on slide 3 onward.
  \end{itemize}

  \onslide<3->{
    \alert<.->{This appears at the same time as the third bullet.}
  }
\end{frame}

\begin{frame}
  \frametitle{More Convoluted Transitions}
  \onslide<1->{Always shown from slide 1 on (\texttt{onslide}).}\\
  \uncover<2->{Uncovered from slide 2 on -- reserves its space even before that (\texttt{uncover}).}\\
  \visible<3->{Visible only from slide 3 on -- like uncover, but never ``ghosted" (\texttt{visible}).}\\
  \invisible<4->{Hidden from slide 4 on -- the opposite of visible (\texttt{invisible}).}\\
  \alt<2>{On slide 2 specifically.}{Not on slide 2.}\quad(\texttt{alt})
\end{frame}

\end{document}
